\section{Basic integration rules}

Basic integration rules come from basic differentiation rules read backward. The derivative of a constant is zero, the derivative of a power follows the power rule, and the derivative of a sum is the sum of derivatives. Integration reverses these facts.

\subsection*{The power rule for integration}

For powers of \(x\), the main rule is:

\begin{theorembox}
If \(p\ne -1\), then
\[
\int x^p\,dx=\frac{x^{p+1}}{p+1}+C.
\]
\end{theorembox}

The condition \(p\ne -1\) is essential. When \(p=-1\), the integrand is \(1/x\), and its integral is logarithmic:
\[
\int \frac1x\,dx=\ln|x|+C.
\]

\begin{examplebox}
Compute
\[
\int x^5\,dx.
\]
Here \(p=5\), so
\[
\int x^5\,dx=\frac{x^6}{6}+C.
\]
Check:
\[
\frac{d}{dx}\left(\frac{x^6}{6}\right)=x^5.
\]
\end{examplebox}

\subsection*{Constants and sums}

Constants can be pulled out of integrals:
\[
\int c f(x)\,dx=c\int f(x)\,dx.
\]
Sums and differences can be integrated term by term:
\[
\int (f(x)+g(x))\,dx=\int f(x)\,dx+\int g(x)\,dx.
\]

\begin{examplebox}
Compute
\[
\int (5x^4-3x^2+7)\,dx.
\]
Integrate term by term:
\[
\int 5x^4\,dx=x^5,
\]
\[
\int (-3x^2)\,dx=-x^3,
\]
\[
\int 7\,dx=7x.
\]
Therefore
\[
\int (5x^4-3x^2+7)\,dx=x^5-x^3+7x+C.
\]
\end{examplebox}

\subsection*{Rewriting before integrating}

Many integrals become simple only after rewriting roots and fractions as powers.

\begin{examplebox}
Compute
\[
\int \frac{x\sqrt{x}}{\sqrt[3]{x}}\,dx.
\]
Rewrite using exponents:
\[
\frac{x\sqrt{x}}{\sqrt[3]{x}}
=\frac{x\cdot x^{1/2}}{x^{1/3}}
=x^{1+1/2-1/3}=x^{7/6}.
\]
Thus
\[
\int \frac{x\sqrt{x}}{\sqrt[3]{x}}\,dx
=\int x^{7/6}\,dx
=\frac{x^{13/6}}{13/6}+C
=\frac6{13}x^{13/6}+C.
\]
\end{examplebox}

The integral was not difficult after the expression was read correctly. The main work was rewriting the integrand into a power.

\subsection*{Elementary integrals}

The following basic antiderivatives should become familiar:

\begin{center}
\begin{tabular}{ll}
\toprule
Integrand & Antiderivative \\
\midrule
\(x^p\), \(p\ne -1\) & \(\dfrac{x^{p+1}}{p+1}+C\) \\
\(1/x\) & \(\ln|x|+C\) \\
\(e^x\) & \(e^x+C\) \\
\(\sin x\) & \(-\cos x+C\) \\
\(\cos x\) & \(\sin x+C\) \\
\(1/\cos^2 x\) & \(\tan x+C\) \\
\(1/(1+x^2)\) & \(\arctan x+C\) \\
\bottomrule
\end{tabular}
\end{center}

A table is useful only after the integrand has been read. Sometimes the integrand must first be rewritten or transformed.

\begin{summarybox}
When applying basic integration rules, ask:
\begin{itemize}
\item Can the integrand be rewritten as powers of \(x\)?
\item Is there a term \(1/x\), which gives a logarithm?
\item Can I integrate term by term?
\item Did I include the constant \(C\)?
\item Can I check the result by differentiating?
\end{itemize}
\end{summarybox}

Basic rules are the foundation. More advanced methods often aim to transform an integral into one of these basic forms.